\section{Vocabulário: Fowler e os dublês}

O vocabulário vem de Martin Fowler
(\citetitle{artigo:testdouble}, \citetitle{artigo:mocks-arent-stubs}) e
de \citeauthor{livro:meszaros2007}.

\begin{itemize}
\item \termo{SUT}{\emph{System Under Test}: o trecho de código que o teste
  está exercitando.} (o objeto sob teste) e o \emph{colaborador} (o que o
  SUT usa).
\item \termo{Fixture}{O estado inicial conhecido de que o teste precisa
  para rodar: objetos criados, dados carregados.}: o estado de partida do
  teste.
\item As quatro fases do teste: \emph{setup}, \emph{exercise},
  \emph{verify}, \emph{teardown}.
\end{itemize}

Os cinco dublês de teste (\emph{test doubles}):

\begin{description}
\item[Dummy] preenche um parâmetro e nunca é de fato usado.
\item[Stub]\termo{Stub}{Dublê que devolve respostas prontas às chamadas
  feitas durante o teste.} devolve respostas prontas.
\item[Fake] tem implementação funcional, mas simplificada (por exemplo um
  repositório em memória).
\item[Spy] registra como foi chamado, para o teste verificar
  \emph{estado} depois.
\item[Mock]\termo{Mock}{Dublê programado com as chamadas que espera
  receber; o teste verifica o \emph{comportamento}.} verifica
  \emph{comportamento}: se as chamadas esperadas aconteceram.
\end{description}

Verificar estado (spy) ou comportamento (mock) é uma escolha de projeto.
Na prática do dojo: objetos reais sempre que possível, e dublês só nas
bordas (banco, rede, relógio).
